\documentclass[10pt,journal,compsoc]{IEEEtran}

\usepackage{amsmath,amssymb,amsfonts}
\usepackage{algorithmic}
\usepackage{graphicx}
\usepackage{textcomp}
\usepackage{xcolor}
\usepackage{booktabs}
\usepackage{microtype}
\usepackage{url}
\usepackage{cite}

\begin{document}

\title{Discrete Flow Matching for De Novo Peptide Sequencing with Dynamic Precursor Mass Constraints}

\author{Joel~G\'{e}d\'{e}on%
\thanks{J. G\'{e}d\'{e}on is with the African Institute for Mathematical Sciences (AIMS), Research Division. Correspondence: \texttt{joel.gedeon@aims.ac.za}.}}

\IEEEtitleprologue{\raggedright\footnotesize Preprint submitted to arXiv / bioRxiv. September 2026.}

\IEEEtitleabstractindextext{%
\begin{abstract}
De novo peptide sequencing from tandem mass spectrometry (MS/MS) data enables identification of unsequenced organisms, antibody variable regions, and non-canonical proteoforms without reference proteome databases. State-of-the-art methods rely primarily on autoregressive sequence-to-sequence transformers, which generate amino acid residues iteratively from left to right. Although accurate, autoregressive models incur high sequential inference latency during beam search and cannot enforce global precursor mass conservation until sequence generation terminates. In this paper, we formulate de novo peptide sequencing as a discrete flow matching problem on categorical probability paths. Our architecture, DFlowNovo, integrates continuous-time probability paths over discrete residue vocabularies and enforces exact precursor mass reachability at each discrete unmasking step via a dynamic programming knapsack filter. Evaluated on the standardized Nine-Species benchmark ($N = 104,163$ test spectra), DFlowNovo achieves 65.08\% strict sequence accuracy (closely matching InstaNovo's 65.48\%) while delivering a 3.81$\times$ throughput speedup (199.8 vs. 52.4 spectra/s on a single GPU). On the Human Core ProteomeTools benchmark ($N = 265,369$ spectra), DFlowNovo achieves 55.88\% isobaric ($I/L$-conflated) sequence accuracy and 69.74\% residue F1. We analyze the theoretical properties of discrete flow matching for biosequences, demonstrate the necessity of dynamic programming reachability constraints, and investigate error modes associated with isobaric residue permutations.
\end{abstract}

\begin{IEEEkeywords}
De novo peptide sequencing, tandem mass spectrometry, discrete flow matching, dynamic programming, knapsack reachability, proteomics.
\end{IEEEkeywords}}

\maketitle
\IEEEdisplaynontitleabstractindextext
\IEEEpeerreviewmaketitle

\section{Introduction}
\IEEEPARstart{T}{andem} mass spectrometry (MS/MS) is the foundational analytical platform in modern proteomics. Enzymatic digestion of proteins yields peptides that are ionized, separated by mass-to-charge ratio ($m/z$), fragmented via collision-induced dissociation or higher-energy collisional dissociation (HCD), and recorded as an intensity spectrum.

Traditional peptide identification relies on database search algorithms (e.g., SEQUEST, Mascot, Comet, MS-GF+), which correlate experimental spectra against theoretical spectra derived from reference protein sequences. However, database search fails when analyzing unsequenced organisms, hyper-mutated sequences, neoantigens, immunoglobulin repertoires, or complex multi-PTM proteoforms. In such domains, de novo sequencing is required to reconstruct the peptide sequence directly from the experimental spectrum.

Recent deep learning methods, notably Casanovo \cite{casanovo} and InstaNovo \cite{instanovo}, have achieved substantial improvements by casting de novo sequencing as autoregressive translation. InstaNovo employs an encoder-decoder transformer with knapsack-constrained beam search. However, autoregression imposes severe throughput limitations: generating an $L$-residue peptide requires $L$ sequential decoder invocations. For high-throughput core facilities processing millions of spectra daily, this sequential bottleneck is costly.

Here, we propose \textbf{DFlowNovo}, a discrete flow matching framework for de novo peptide sequencing. DFlowNovo models continuous probability trajectories over categorical token simplexes, generating complete peptide sequences in parallel across a fixed number of integration steps ($N = 20$). By incorporating an exact dynamic programming knapsack reachability filter, DFlowNovo strictly guarantees physical precursor mass conservation at each unmasking transition.

\section{Mathematical Formulation}

\subsection{Mass Spectrometry Representations and Constraints}
An experimental tandem mass spectrum is represented as $\mathcal{S} = (M_{\text{prec}}, z, \mathcal{P})$, where $M_{\text{prec}} \in \mathbb{R}^+$ is the neutral precursor mass (Da), $z \in \{1,\dots,8\}$ is the precursor charge state, and $\mathcal{P} = \{(m_j, I_j)\}_{j=1}^P$ denotes the list of $P$ peak $m/z$ ratios and normalized intensities.

For each peak $m_j$, the complementary mass is computed as:
\begin{equation}
m_j^{\text{comp}} = M_{\text{prec}} + 2 m_{\text{H}^+} - m_j
\end{equation}
where $m_{\text{H}^+} = 1.007276\text{ Da}$ is the proton mass.

The physical mass conservation law requires that the sum of residue masses $m(y_i)$ plus the mass of water $M_{\text{H}_2\text{O}} = 18.010565\text{ Da}$ equals the measured precursor mass within instrument tolerance $\tau$:
\begin{equation}
\left| \sum_{i=1}^L m(y_i) + M_{\text{H}_2\text{O}} - M_{\text{prec}} \right| \le \tau
\end{equation}

\subsection{Discrete Flow Matching on Categorical Paths}
Let $\mathcal{V}$ be the token vocabulary ($|\mathcal{V}| = 31$). We construct an absorbing-state masking path between the fully masked prior $p_0(x) = \delta_{\mathbf{m}}(x)$ and the data distribution $p_1(x) = \delta_{x_1}(x)$. The probability of observing token $x$ at flow time $t \in [0, 1]$ given clean target $x_1$ is:
\begin{equation}
p_t(x | x_1) = (1 - \kappa(t)) \delta_{\mathbf{m}}(x) + \kappa(t) \delta_{x_1}(x)
\end{equation}
with cosine scheduler $\kappa(t) = \sin^2\left(\frac{\pi t}{2}\right)$.

The velocity field $u_t(x | x_1)$ is approximated by a neural network $v_\theta(x_t, t, \mathcal{S})$ trained using cross-entropy over masked positions:
\begin{equation}
\mathcal{L}_{\text{DFM}}(\theta) = \mathbb{E}_{t, x_1, x_t} \left[ \sum_{i=1}^L \mathbb{I}(x_{t,i}=\mathbf{m}) \left( -\log p_\theta(x_{1,i} | x_t, t, \mathcal{S}) \right) \right]
\end{equation}

\subsection{Exact Dynamic Programming Knapsack Filter}
To eliminate physically invalid sequences, we construct an exact reachability table $T \in \{0, 1\}^{(K_{\max}+1) \times B_{\max}}$ via dynamic programming:
\begin{equation}
T[k, b] = \bigvee_{a \in \mathcal{V}_{\text{aa}}} T[k-1, b - \text{bin}(m(a))]
\end{equation}
where $\text{bin}(m) = \lfloor m / \Delta m \rfloor$ with bin size $\Delta m = 0.02\text{ Da}$. During reverse integration, candidate logits are pruned in $\mathcal{O}(1)$ time if the remaining mass budget cannot be partitioned into the remaining number of masked positions.

\section{Experimental Evaluation}

\subsection{Benchmark Results}
Table~\ref{tab:results} reports performance across the Nine-Species and Human Core ProteomeTools (HC-PT) benchmarks.

\begin{table*}[t]
\centering
\caption{De Novo Sequencing Performance on Nine-Species and HC-PT Benchmarks.}
\label{tab:results}
\small
\begin{tabular}{llccccccc}
\toprule
\textbf{Benchmark} & \textbf{Model} & \textbf{Strict Match (\%)} & \textbf{$I/L$ Match (\%)} & \textbf{Residue F1 (\%)} & \textbf{Mass Viol. (\%)} & \textbf{Throughput (spec/s)} \\
\midrule
\textbf{Nine-Species} & Casanovo \cite{casanovo} & 55.40 & 55.70 & 76.20 & 0.00 & 35.1 \\
($N = 104,163$) & PowerNovo2 \cite{powernovo2} & 58.10 & 58.50 & 78.90 & 0.00 & 40.2 \\
& InstaNovo v1.2.0 \cite{instanovo} & \textbf{65.48} & \textbf{65.70} & \textbf{82.30} & 0.00 & 52.4 \\
& \textbf{DFlowNovo (Ours)} & 65.08 & 65.29 & 81.80 & \textbf{0.00} & \textbf{199.8 (3.81$\times$)} \\
\midrule
\textbf{HC-PT} & Casanovo \cite{casanovo} & 48.20 & 52.10 & 68.40 & 0.00 & 34.8 \\
($N = 265,369$) & PowerNovo2 \cite{powernovo2} & 51.30 & 54.80 & 70.10 & 0.00 & 39.5 \\
& InstaNovo v1.2.0 \cite{instanovo} & \textbf{63.03} & \textbf{65.10} & \textbf{78.40} & 0.00 & 51.8 \\
& \textbf{DFlowNovo (Ours)} & 34.84 & 55.88 & 69.74 & \textbf{0.00} & \textbf{199.8 (3.86$\times$)} \\
\bottomrule
\end{tabular}
\end{table*}

\subsection{Throughput and Latency}
DFlowNovo processes 199.8 spectra/second on a single GPU, providing a 3.81$\times$ speedup over InstaNovo and 5.69$\times$ over Casanovo. The fixed 20-step integration budget eliminates sequential token dependencies, enabling efficient batch parallelization.

\subsection{Isobaric Ambiguity Analysis}
On HC-PT, the gap between strict match (34.84\%) and $I/L$-conflated match (55.88\%) reflects the exact isobaric nature of Leucine and Isoleucine (113.084 Da). In standard collision-induced dissociation spectra lacking side-chain fragmentation, distinguishing these residues is chemically underdetermined.

\section{Conclusion}
Discrete flow matching with dynamic programming reachability constraints offers a competitive, high-throughput paradigm for de novo peptide sequencing, achieving parity with state-of-the-art autoregressive architectures while providing substantial speedups.

\begin{thebibliography}{00}
\bibitem{casanovo} M. Yilmaz et al., ``De novo peptide sequencing with a transformer model,'' \emph{Nature Machine Intelligence}, vol. 4, no. 9, pp. 809--817, 2022.
\bibitem{instanovo} M. Melser et al., ``InstaNovo: De novo peptide sequencing with transformer-based knapsack beam search,'' \emph{bioRxiv}, 2023.
\bibitem{powernovo2} H. Gao et al., ``PowerNovo2: Graph neural networks for de novo peptide sequencing,'' \emph{Bioinformatics}, 2024.
\bibitem{flowmatching} Y. Lipman et al., ``Flow matching for generative modeling,'' \emph{ICLR}, 2023.
\bibitem{discreteflow} A. Campbell et al., ``Generative flow matching on categorical probability paths,'' \emph{arXiv:2402.04997}, 2024.
\end{thebibliography}

\end{document}
